\section{Certificate-degree proofs}\label{app:math}

This appendix proves the mathematical statements used in the case study. All polynomial rings are over $\C$ unless rational coefficients are specified. Tensor rank is the minimum number of completely decomposable summands. The equations and degree convention are those of \cref{eq:raw-cp,eq:unit-certificate}; the order $n$ and proposed number of summands $R$ are separate parameters.

\subsection{The inherited universal lower bound}

\begin{lemma}[Distinct-label functional]\label{lem:lower}
For every order-$n$ tensor $T$ and integer $R\ge1$, the original CP-$R$ system satisfies
\[
                         \delta_R(T)\ge n(R+1).
\]
The statement is independent of the finite mode dimensions. When $\rank(T)\le R$, the inequality holds with $\delta_R(T)=\infty$ by the convention in \cref{sec:case}.
\end{lemma}

\begin{proof}
On monomials of total degree at most $n(R+1)-1$, define a linear functional $L$ as follows. A monomial is admissible if, at every occupied summand label $r$, it contains exactly one variable from each of the $n$ modes. Thus each occupied label determines an output multi-index $J_r$. For $t$ occupied labels, set
\[
 L\!\left(\prod_{r\text{ occupied}}\prod_{i=1}^n x_{r,i,j_{r,i}}\right)
    =\frac{\prod_{r\text{ occupied}}T_{J_r}}{(R)_t},
 \qquad (R)_t=R(R-1)\cdots(R-t+1).
\]
Set $L(1)=1$ and assign zero to every nonadmissible monomial. There are at most $R$ occupied labels, so all denominators used are nonzero.

Let $M$ be any multiplier monomial of degree at most $nR-1$. If $M$ is admissible, it occupies $t\le R-1$ labels. In the product $MQ_J$, the $t$ summands using an already occupied label are nonadmissible. Each of the remaining $R-t$ labels contributes
\[
 \frac{T_J\prod_{r\text{ occupied}}T_{J_r}}{(R)_{t+1}}.
\]
Their sum is $T_JL(M)$. If $M$ is nonadmissible, multiplication by one variable from every mode at a single label cannot make it admissible: a partially occupied label acquires a repeated mode, an already repeated mode persists, and any other bad label is unchanged. Hence both $L(MQ_J)$ and $T_JL(M)$ are zero. Thus $L(MF_J)=0$ for every permitted $M$ and $J$.

Every nonzero $F_J$ has degree $n$, so a certificate of product degree at most $n(R+1)-1$ would have multipliers of degree at most $nR-1$. Applying $L$ to that identity would give $1=0$. This proves the bound.
\end{proof}

This argument is the lower functional inherited from the atlas. It is written here in full because its independence from mode dimensions is needed for both the pencil family and contractions with passive modes.

\subsection{The flattening certificate and balanced multipliers}

\begin{lemma}[Minor certificate]\label{lem:minor}
If a bipartition flattening of $T$ has a nonzero $(R+1)\times(R+1)$ minor, then $\delta_R(T)\le n(R+1)$.
\end{lemma}

\begin{proof}
Write $q=R+1$ and select the corresponding submatrices $Q_0(x)$ and $T_0$. The flattening of $Q(x)$ is a sum of $R$ rank-one matrices, so $\det Q_0(x)$ is identically zero. Replace the rows of $T_0$ by rows of $Q_0$ one at a time. Multilinearity of the determinant gives
\[
 -\det T_0=\det Q_0-\det T_0=\sum_J c_J(x)F_J(x),
\]
where each cofactor $c_J$ uses at most $R$ entries of $Q$ and remaining constant entries of $T$. Therefore $\deg c_J\le nR$. Dividing by the nonzero constant $-\det T_0$ gives the required certificate.
\end{proof}

\begin{lemma}[Balance reduction]\label{lem:balance}
If an original CP-$R$ system admits a unit certificate of degree at most $D$, it admits one at the same degree whose multiplier monomials have equal degrees in all modes separately at each summand label. The minimum finite certificate degree is consequently a multiple of $n$.
\end{lemma}

\begin{proof}
For each label $r$, act on its mode vectors by scalars $\lambda_{r,1},\ldots,\lambda_{r,n}$ with $\prod_i\lambda_{r,i}=1$. Every $F_J$ and the constant one are invariant. Projecting a polynomial identity to the zero character of this torus projects each multiplier to its invariant monomials and does not increase degree. A monomial is invariant precisely when its mode degrees agree within each label. Such a monomial has degree $n\sum_r t_r$ for nonnegative integers $t_r$. A nonzero balanced multiplier times $F_J$ has degree $n(1+\max\sum_r t_r)$, a multiple of $n$. Applying the projection to a minimum-degree certificate proves the last statement. Averaging further over permutations of the $R$ labels is also valid in characteristic zero and gives the orbit representation used by the computation.
\end{proof}

\subsection{An explicit boundary certificate}

Consider the tensor $S\in\C^2\otimes\C^2\otimes\C^2$ with slices $I_2,J_2$. With indices $\alpha,i,j\in\{0,1\}$, let
\[
 F_{\alpha i j}=\sum_{r=1}^{2}a_{r,\alpha}b_{r,i}c_{r,j}-S_{\alpha i j}.
\]
For a monomial $p$ in the six symbols $a_0,a_1,b_0,b_1,c_0,c_1$, write $p^{(r)}$ for its evaluation at label $r$, and define
\[
                  A(p,q)=\frac{p^{(1)}q^{(2)}+p^{(2)}q^{(1)}}{2}.
\]
This convention also applies when $p=q$.

\begin{lemma}[Degree-9 boundary identity]\label{lem:base}
The sum of $\kappa A(p,q)F_{\alpha i j}$ over the rows of \cref{tab:base-certificate} equals one. Every multiplier is label-balanced and has degree at most six. Thus $\delta_2(S)=9$.
\end{lemma}

\begin{table}[htbp]
\centering\small
\begin{tabular}{@{}rlll@{}}
\toprule
$\kappa$ & $p$ & $q$ & $(\alpha,i,j)$ \\
\midrule
$-1$ & $1$ & $1$ & $(1,0,1)$ \\
$4$ & $1$ & $a_1b_1c_1$ & $(0,0,1)$ \\
$-2$ & $1$ & $a_1b_0c_1$ & $(0,1,1)$ \\
$2$ & $1$ & $a_0a_1b_0b_1c_1^2$ & $(1,0,1)$ \\
$-2$ & $1$ & $a_0a_1b_0^2c_1^2$ & $(0,1,0)$ \\
$2$ & $1$ & $a_0a_1b_0^2c_1^2$ & $(1,0,0)$ \\
$-2$ & $1$ & $a_0a_1b_0^2c_1^2$ & $(1,1,1)$ \\
$2$ & $1$ & $a_0a_1b_0^2c_0c_1$ & $(0,1,1)$ \\
$-2$ & $1$ & $a_0a_1b_0^2c_0c_1$ & $(1,0,1)$ \\
$4$ & $1$ & $a_0^2b_0^2c_1^2$ & $(1,1,0)$ \\
$-4$ & $1$ & $a_0^2b_0^2c_0c_1$ & $(1,1,1)$ \\
$2$ & $a_1b_1c_1$ & $a_1b_0c_1$ & $(0,0,1)$ \\
$4$ & $a_1b_1c_1$ & $a_0b_0c_1$ & $(0,0,0)$ \\
$-4$ & $a_1b_1c_0$ & $a_0b_0c_1$ & $(0,0,1)$ \\
$2$ & $a_1b_0c_1$ & $a_1b_0c_1$ & $(0,0,0)$ \\
$-2$ & $a_1b_0c_1$ & $a_1b_0c_1$ & $(0,1,1)$ \\
$-2$ & $a_1b_0c_1$ & $a_1b_0c_0$ & $(0,0,1)$ \\
$-2$ & $a_1b_0c_1$ & $a_0b_1c_1$ & $(0,0,0)$ \\
$2$ & $a_1b_0c_1$ & $a_0b_1c_0$ & $(0,0,1)$ \\
\bottomrule
\end{tabular}
\caption{The complete boundary certificate. Each row contributes $\kappa A(p,q)F_{\alpha i j}$; their sum is the constant one.}
\label{tab:base-certificate}
\end{table}

\begin{proof}
Substitute $S_{0ij}=\mathbf1_{i=j}$ and $S_{1ij}=\mathbf1_{(i,j)=(0,1)}$ in the displayed sum and expand. All nonconstant coefficients cancel and the constant coefficient is one. The table specifies the entire rational identity, so this verification requires only polynomial multiplication and collection; the companion also implements that expansion directly. The mode degrees in each displayed $p$ and $q$ agree, and each $A(p,q)$ has degree at most six. The product degree is at most nine. A certificate proves inconsistency, and \cref{lem:lower} supplies the matching lower bound.
\end{proof}

Swapping $e_0$ and $e_1$ in the third mode sends $S$ to $W_3$. This is a constant linear change of variables, so the same certificate, with relabeled coordinates and equations, applies to $W_3$ with the same degrees and balance.

\subsection{A degree-controlled transfer lemma}

\begin{lemma}[Balanced transfer]\label{lem:transfer}
Let a $k$-mode tensor $S$ have a CP-$R$ unit certificate whose multipliers are label-balanced and have degree at most $k(q-1)$. Let $n\ge k$, and suppose an $n$-mode tensor $T$ maps to $S$ by constant linear maps on $k$ retained modes and constant linear functionals on the remaining modes. Then the original CP-$R$ system for $T$ has a unit certificate of product degree at most $nq$.
\end{lemma}

\begin{proof}
Permute modes so that the retained modes are $1,\ldots,k$. Write $A_i:V_i\to U_i$ for their linear maps and $\ell_i:V_i\to\C$ for the other functionals. For each summand label, substitute into the base parameters
\[
 y_{r,1}=A_1x_{r,1}\prod_{i=k+1}^{n}\ell_i(x_{r,i}),
 \qquad y_{r,j}=A_jx_{r,j}\quad(2\le j\le k).
\]
The first vector has degree at most $n-k+1$ and each other vector has degree at most one. A balanced multiplier monomial with degree $t_r$ in each of the $k$ modes at label $r$ becomes a polynomial of degree at most
\[
 \sum_r\bigl((n-k+1)t_r+(k-1)t_r\bigr)=n\sum_r t_r.
\]
Hence every substituted multiplier has degree at most $n(q-1)$.

Let $L=(\bigotimes_{i=1}^k A_i)\otimes(\bigotimes_{i=k+1}^n\ell_i)$ be the resulting constant linear map on tensor entries. By construction, $Q^{S}(y)=L(Q^{T}(x))$ and $S=L(T)$. Therefore each substituted base equation is a constant linear combination of the original equations:
\[
 F^S_K(y(x))=\sum_J L_{KJ}F^T_J(x).
\]
Substitution in the base unit identity gives
\[
 1=\sum_K h_K(y(x))\sum_J L_{KJ}F^T_J(x).
\]
Each multiplier has degree at most $n(q-1)$ and each original equation degree $n$, proving the bound. The construction is a polynomial identity even when some contracted parameter forms vanish.
\end{proof}

\begin{corollary}[Normalized certificate degree under contraction]\label{cor:degree-transfer}
Let $R\ge1$ and $n\ge k\ge1$. Let $T$ be an $n$-mode tensor and $S$ a $k$-mode tensor over $\C$, with $S$ obtained from $T$ by constant linear maps on $k$ retained modes and constant linear functionals on the other modes. If $\delta_R(S)<\infty$, then
\begin{equation}\label{eq:degree-transfer}
 \delta_R(T)\le\frac{n}{k}\,\delta_R(S),
 \qquad\text{equivalently}\qquad
 \frac{\delta_R(T)}{n}\le\frac{\delta_R(S)}{k}.
\end{equation}
Both degrees are measured in their respective original CP-$R$ parameter systems.
\end{corollary}

\begin{proof}
Put $D=\delta_R(S)$ and choose a certificate attaining this finite minimum. Applied to the $k$-mode system, \cref{lem:balance} gives a balanced certificate of degree at most $D$ and shows that $D=kq$ for an integer $q$. Its multipliers have degree at most $D-k=k(q-1)$, because every original equation has degree $k$. \Cref{lem:transfer} therefore supplies a certificate for $T$ of degree at most $nq=(n/k)D$. In particular the transferred system is inconsistent as well. No separate hypothesis that the chosen certificate is already balanced is needed.
\end{proof}

\begin{corollary}\label{cor:wn}
For every $n\ge3$, the tensor $W_n$ in \cref{eq:wn} satisfies $\delta_2(W_n)=3n$.
\end{corollary}

\begin{proof}
Contract modes $4,\ldots,n$ by $e_0^*$. The result is $W_3$. By \cref{lem:base}, $\delta_2(W_3)=9$. Applying \cref{cor:degree-transfer} with $k=3$ and $R=2$ gives an upper certificate of degree $(n/3)\cdot9=3n$. The certificate proves the fiber is empty, and \cref{lem:lower} gives the opposite inequality. The flattening formula and degeneration in \cref{sec:boundary} show that every flattening has rank at most two and the border rank is two. No secant classification is needed for this corollary.
\end{proof}

\subsection{Proof of the full two-term theorem}

\begin{proof}[Proof of \cref{thm:main}]
The lower bound $3n$ is \cref{lem:lower} with $R=2$. If a flattening has rank greater than two, \cref{lem:minor} gives the matching upper bound. This covers $n=2$: a matrix of rank greater than two has a nonzero $3\times3$ minor, so the two lemmas give $\delta_2(T)=6=3n$. We may therefore assume $n\ge3$. It remains to consider tensors for which every bipartition flattening has rank at most two.

The set-theoretic flattening theorem of Landsberg and Manivel \cite[arXiv version, Theorem~5.1]{landsbergmanivel} places such a tensor in the second secant variety of the Segre variety. Raicu proves the stronger ideal-theoretic statement \cite[Theorem~4.1 and Corollary~4.2]{raicu}. Write $X$ for this Segre variety. Points of $\sigma_2(X)\setminus\tau(X)$ have rank two \cite[arXiv version, Corollary~1]{ballicobernardi}. Hence $\rank(T)>2$ forces $[T]\in\tau(X)$. By \cite[arXiv version, Notation~1]{ballicobernardi}, $\tau(X)$ is the union of tangent spaces, so $T$ has the form \cref{eq:tangent-main} with nonzero $v_i$ and appropriate $u_i$.

Absorb components of $u_i$ parallel to $v_i$ into $\lambda$. Inactive modes now have $u_i=0$. There must be at least three active modes: with at most two, factoring out passive vectors leaves a vector or a matrix supported on a two-dimensional space in each active mode, of rank at most two. Select three active modes. On each, choose a linear map sending $v_i$ to $e_0$ and $u_i$ to $e_1$. On every unselected mode, choose a functional taking $v_i$ to one and $u_i$ to zero. The resulting three-mode tensor is
\[
                 \lambda e_0\otimes e_0\otimes e_0+W_3.
\]
In the first retained mode, the linear shear fixing $e_0$ and sending $e_1$ to $e_1-\lambda e_0$ removes the pure term. The target therefore maps to $W_3$ by maps of the form required in \cref{lem:transfer}. With the degree-9 certificate of \cref{lem:base}, this gives the upper bound $3n$.

Dimension-one factors may be suppressed when applying the classical geometric statement, since they do not change tensor rank or flattening rank. They remain among the original parameter modes and are contracted in the transfer, so each is still charged in the factor $n$. If fewer than three nonscalar active modes remain in this case, the rank would be at most two, already excluded. This also accounts for the endpoint dimensions.
\end{proof}

\subsection{Uniform certificates for nilpotent pencils}\label{app:pencil}

We prove \cref{thm:pencil} constructively. For a polynomial $p$, write $p\in\mathcal I_{\le d}$ if it has an expression as a linear combination of the original coefficient equations in which every multiplier--equation product has degree at most $d$. Sums preserve the larger bound, and multiplication by a degree-$e$ polynomial raises it by at most $e$.

First choose a constant matrix $H$ satisfying $NH-HN=N$. In Jordan form, take diagonal entries $1,2,\ldots$ successively within each block and then conjugate $H$ back. Introduce square matrices $B,C$ with columns $b_r,c_r$, and diagonal matrices
\[
 A=\diag(a_1,\ldots,a_m),\qquad D=\diag(d_1,\ldots,d_m).
\]
Set
\[
 X=BAC^{\mathsf T},\quad Y=BDC^{\mathsf T},\quad
 F=X-I_m,\quad G=Y-N,
\]
and
\[
 \beta=\det B,\quad \gamma=\det C,\quad U=\beta\gamma,
 \quad P=\prod_{r=1}^m a_r.
\]
The entries of $F,G$ are exactly the original cubic generators. Adjugate identities give
\begin{align*}
 L&=(NF-G)\adj(C^{\mathsf T})=\gamma(NBA-BD),\\
 R'&=\adj(B)(FN-G)=\beta(AC^{\mathsf T}N-DC^{\mathsf T}).
\end{align*}
Their entries belong to $\mathcal I_{\le m+2}$. Write $L_r$ for column $r$ of $L$ and $R'_r$ for row $r$ of $R'$. Then
\begin{align*}
 S_r&=\gamma R'_rHb_r-\beta c_r^{\mathsf T}HL_r
      =Ua_r c_r^{\mathsf T}(NH-HN)b_r
      =Ua_r c_r^{\mathsf T}Nb_r,\\
 T_r&=\beta c_r^{\mathsf T}L_r
      =U\bigl(a_r c_r^{\mathsf T}Nb_r-d_r c_r^{\mathsf T}b_r\bigr).
\end{align*}
Both belong to $\mathcal I_{\le2m+3}$. Setting $z_r=c_r^{\mathsf T}b_r$ gives
\[
 W_r=S_r-T_r=Ud_rz_r\in\mathcal I_{\le2m+3}.
\]
Also
\[
 V_r=(\adj(B)FB)_{rr}=\beta(a_rz_r-1)
       \in\mathcal I_{\le m+3}.
\]
Consequently,
\[
 K_r=a_rW_r-d_r\gamma V_r=Ud_r\in\mathcal I_{\le2m+4}.
\]
Choose fixed indices $i,j$ with $N_{ij}\ne0$. The identity
\[
 Z=\frac{\sum_r(b_r)_i(c_r)_jK_r-UG_{ij}}{N_{ij}}=U
\]
places $U$ in $\mathcal I_{\le2m+6}$. The denominator is a nonzero field constant. Multiplying by $P$ yields $PZ=\det X\in\mathcal I_{\le3m+6}$. Finally, determinant telescoping gives $\det X-1\in\mathcal I_{\le3m}$. Subtraction produces the unit certificate
\[
                      1=PZ-(\det X-1)\in\mathcal I_{\le3m+6}.
\]
Every displayed intermediate expression has an explicit combination of the original equations, so substitution gives all final multipliers, including at singular parameter matrices.

\begin{table}[htbp]
\centering\small
\begin{tabular}{@{}ll@{}}
\toprule
Derived expression & Maximum original product degree \\
\midrule
Entries of $L,R'$ & $m+2$ \\
$S_r,T_r,W_r$ & $2m+3$ \\
$V_r$ & $m+3$ \\
$K_r=Ud_r$ & $2m+4$ \\
$Z=U$ & $2m+6$ \\
$PZ=\det X$ & $3m+6$ \\
$\det X-1$ & $3m$ \\
Unit identity & $3m+6$ \\
\bottomrule
\end{tabular}
\caption{Degree ledger for the nilpotent-pencil certificate. Derived equations are charged through their representations in the original cubic generators.}
\label{tab:degree-ledger}
\end{table}

For completeness, $(I_m,N)$ has no rank-$m$ decomposition. If $I_m=BAC^{\mathsf T}$, all three square factors are invertible. Then $N=BDC^{\mathsf T}=BDA^{-1}B^{-1}$ is diagonalizable, impossible for a nonzero nilpotent matrix. In Jordan coordinates, the perturbations $N+\varepsilon\diag(1,\ldots,m)$ have distinct eigenvalues for $\varepsilon\ne0$; the corresponding pencils have rank $m$ and converge to $(I_m,N)$. The identity slice supplies the lower border-rank bound $m$, and the three flattening ranks are $2,m,m$.

Apply \cref{lem:lower} with order three and $R=m$ to get the lower degree $3m+3$. \Cref{lem:balance} makes the minimum degree divisible by three. Together with the constructive upper bound, this proves \cref{thm:pencil}. The rational certificates at $m=2,3,4$ select the lower endpoint for single Jordan blocks; the uniform theorem retains both possibilities.
